\documentclass[
a4paper, % Stock and paper size.
12pt, % Type size.
article,
% oneside, 
onecolumn, % Only one column of text on a page.
% openright, % Each chapter will start on a recto page.
% openleft, % Each chapter will start on a verso page.
openany, % A chapter may start on either a recto or verso page.
]{memoir}

%%% PACKAGES 
%%%------------------------------------------------------------------------------

\usepackage[utf8]{inputenc} % If utf8 encoding
% \usepackage[lantin1]{inputenc} % If not utf8 encoding, then this is probably the way to go
\usepackage[T1]{fontenc}    %
\usepackage[english,russian]{babel} % English please
\usepackage[final]{microtype} % Less badboxes

% \usepackage{kpfonts} %Font

\usepackage{amsmath,amssymb,mathtools} % Math

%%% Figures and colors
\usepackage[usenames,dvipsnames,svgnames,table,rgb]{xcolor}
\usepackage{tikz} % Figures
\usepackage{graphicx}  % Include figures
\usepackage{wrapfig}
\graphicspath{{img/}{../img/}}

%%% PAGE LAYOUT 
%%%------------------------------------------------------------------------------

\setlrmarginsandblock{0.08\paperwidth}{*}{1} % Left and right margin
\setulmarginsandblock{0.15\paperwidth}{*}{1}  % Upper and lower margin
\checkandfixthelayout
%%% indenting
\setlength{\parindent}{0em}
\setlength{\parskip}{0.5em}

%%% INTERNAL HYPERLINKS
%%%-------------------------------------------------------------------------------

\usepackage{hyperref}   % Internal hyperlinks
\hypersetup{
pdfborder={0 0 0},      % No borders around internal hyperlinks
pdfauthor={I am the Author} % author
}
% ----------- hyper ref -----------
\usepackage{hyperref}


\newcommand{\linkcolor}{blue}
\newcommand{\citecolor}{blue}
\newcommand{\filecolor}{magenta}
\newcommand{\urlcolor}{NavyBlue}
\hypersetup{				% Гиперссылки
	unicode=true,           % русские буквы в раздела PDF\\
	pdfstartview=FitH,
	colorlinks=true,  % false: ссылки в рамках; true: цветные ссылки
	linkcolor=\linkcolor,         % внутренние ссылки
	citecolor=\citecolor,        % на библиографию
	filecolor=\filecolor,      % на файлы
	urlcolor=\urlcolor,      % на URL
}
\let\oldhref\href
\renewcommand{\href}[2]{\oldhref{#1}{\underline{#2}}}

\usepackage{memhfixc}   %



% Defining a new coordinate system for the page:
%
% --------------------------
% |(-1,1)    (0,1)    (1,1)|
% |                        |
% |(-1,0)    (0,0)    (1,0)|
% |                        |
% |(-1,-1)   (0,-1)  (1,-1)|
% --------------------------
\makeatletter
\def\parsecomma#1,#2\endparsecomma{\def\page@x{#1}\def\page@y{#2}}
\tikzdeclarecoordinatesystem{page}{
    \parsecomma#1\endparsecomma
    \pgfpointanchor{current page}{north east}
    % Save the upper right corner
    \pgf@xc=\pgf@x%
    \pgf@yc=\pgf@y%
    % save the lower left corner
    \pgfpointanchor{current page}{south west}
    \pgf@xb=\pgf@x%
    \pgf@yb=\pgf@y%
    % Transform to the correct placement
    \pgfmathparse{(\pgf@xc-\pgf@xb)/2.*\page@x+(\pgf@xc+\pgf@xb)/2.}
    \expandafter\pgf@x\expandafter=\pgfmathresult pt
    \pgfmathparse{(\pgf@yc-\pgf@yb)/2.*\page@y+(\pgf@yc+\pgf@yb)/2.}
    \expandafter\pgf@y\expandafter=\pgfmathresult pt
}
\makeatother

%%% for beautiful icons
\usepackage{fontawesome}


\begin{document}

\tikz[remember picture,overlay] \node[opacity=0.9,inner sep=0pt] at (page cs:0.8,0.8){\includegraphics[width=0.1\paperwidth]{IshaLogo}};

\section*{Инструкции к первой домашней работе}
\label{sec:homework}
\addcontentsline{toc}{section}{\nameref{sec:homework}}

\subsection*{Партнерская система}
\begin{enumerate}
\item группа в телеграмме будет создана примерно для 10-15 волонтеров йога вира.
\item ваш йога митра (координатор) назначит вам партнера, с которым вы будете работать в
паре.
\item в качестве домашнего задания, вы проведёте практическую сессию вместе со своим
партнером.
\item после того, как вы оба завершите проведение своей сессии, уведомите об этом йога митру.
\end{enumerate}

\subsection*{Структура}

\begin{enumerate}
\item Продолжительность сессии 20 минут.
\item Заранее подготовьте \textbf{любой из видеороликов}, которые вы будете использовать во время демонстрации. Пожалуйста, ознакомьтесь с процессом трансляции видео перед сессией. Воспользуйтесь обучающими видео из документа  <<Набор Йога Виры>>.
\item Представьтесь в начале сессии, так как это описано в сценарии, который у вас есть.
(пожалуйста, максимально придерживайтесь текста сценария)
\item Воспроизведите приготовленное видео  (вы можете остановить видео после 30 секунды).
\item Расскажите о программе «Внутренняя инженерия» так, как это описано в сценарии.
\item Завершите сессию так, как это описано в сценарии.

    \textbf{Вам не нужно рассказывать о своем опыте <<Внутренней инженерии>> на этом занятии. Мы потренируем эту часть отдельно.}
\end{enumerate}


\subsection*{Вы ведете практическую сессию}
\begin{itemize}
\item Проводите репетицию практической сессии так, как если бы вы проводили настоящую.
сессию. Не останавливайте и не прерывайте сессию, даже если вы совершили ошибку —
попробуйте найти решение как можно продолжить дальше \faSmileO.
\item Как только практическая сессия закончена, пожалуйста, уделите минуту, чтобы ответить на следующие вопросы и поделитесь ответами со своим партнером (вы отвечаете на эти вопросы сами, а не партнер).
\begin{enumerate}
\item Какие у вас впечатления после проведения практической сессии? Что происходило
внутри вас?
\item Какие, по-вашему мнению, части сессии могли бы быть сделаны лучше?
\end{enumerate} 
\end{itemize}

\subsection*{Ваш партнер проводит сессию}
\begin{itemize}
\item Принимайте участие словно вы пришли на настоящую сессию в качестве участника. Не
перебивайте его в процессе занятия.
\item Создайте атмосферу доверия, чтобы ваш партнёр смог поделиться своими
впечатлениями и ощущениями без страха осуждения \faSmileO.
\item Если вы заметили что-то, что можно улучшить, убедитесь, что вы сами привнесете это в
свою сессию, когда будете ее проводить.
\item Поблагодарите и поддержите своего партнера после того, как он поделиться своими
ощущениями от проведённой сессии.
\end{itemize} 

\tikz[remember picture,overlay] \node[opacity=0.15,inner sep=0pt] at (current page.center){\includegraphics[width=0.2\paperwidth]{ishaLogo}};
\end{document}

